\chapter{Thực nghiệm và Đánh giá}
\label{ch:evaluation}

Chương này trình bày kết quả trên 200 chương trình vòng lặp Python. Các phương pháp được đối chiếu theo cùng một khẳng định gốc và một nhãn chuẩn cho mỗi chương trình.

\section{Thiết lập đánh giá}
\label{sec:eval_setup}

\subsection{Dữ liệu và nhãn chuẩn}
\label{sec:datasets}

Tập thực nghiệm được lấy từ \texttt{bench\_v1}, bộ dữ liệu gồm 1.000 chương trình. Tập con có 100 mẫu mang nhãn \texttt{VERIFIED} và 100 mẫu mang nhãn \texttt{COUNTEREXAMPLE}. Cả 200 mẫu đều chứa vòng lặp và thuộc họ Code2Inv. Trong đó có 7 chương trình từ kho gốc và 193 biến thể sinh tự động. Danh sách định danh được dùng chung cho các phương pháp.

Nhãn của bộ dữ liệu dựa trên thử nghiệm bằng CPython và, ở các trường hợp phù hợp, kiểm tra ký hiệu bằng Z3/Spacer. Cả 100 mẫu mang nhãn \texttt{COUNTEREXAMPLE} đều có đầu vào phản ví dụ trong manifest.

\subsection{Cấu hình và quy tắc chấm}

Các lượt dùng LLM trong bảng sử dụng MiniCPM-V-4.7-35B-A3B qua AMD Token Lab. Proposed Offline, J, K và L chạy cục bộ. Cả hai chế độ Proposed đều được chấm trên câu lệnh \texttt{assert} gốc. Các khẳng định bổ sung do MiniCPM đề xuất được thống kê riêng.

Một mẫu được tính đúng nếu phán quyết khớp chính xác với nhãn của khẳng định gốc. Gọi $N$ là số mẫu, $N_c$ là số phán quyết đúng và $N_d$ là số mẫu nhận được \texttt{VERIFIED} hoặc \texttt{COUNTEREXAMPLE}. Độ chính xác là $N_c/N$ và độ phủ phán quyết là $N_d/N$. Các trạng thái không kết luận, lỗi dịch và lỗi chạy đều nằm trong mẫu số $N=200$ và được tính vào số mẫu chưa có phán quyết.

Mỗi định danh mẫu đóng góp một kết quả cuối cùng cho mỗi phương pháp.

\section{So sánh trên 200 khẳng định gốc}
\label{sec:overall_benchmark_results}

\begin{table}[htbp]
\centering
\footnotesize
\setlength{\tabcolsep}{3pt}
\caption{Kết quả trên 200 chương trình vòng lặp Python}
\label{tab:benchmark-overall-expanded}
\begin{tabular}{@{}p{0.335\textwidth}rrrrrrrr@{}}
\toprule
\textbf{Phương pháp} & \textbf{Đúng} & \textbf{Sai} & \makecell[r]{\textbf{Chưa}\\\textbf{kết luận}} & \makecell[r]{\textbf{Độ}\\\textbf{chính xác}} & \textbf{Độ phủ} & \textbf{Precision} & \textbf{Recall} & \textbf{F1} \\
\midrule
Proposed + MiniCPM & 200 & 0 & 0 & 100\% & 100\% & 100\% & 100\% & 100\% \\
Proposed Offline & 200 & 0 & 0 & 100\% & 100\% & 100\% & 100\% & 100\% \\
A -- SatLM Direct & 106 & 43 & 51 & 53\% & 74,5\% & 73,02\% & 46\% & 56,44\% \\
B -- Execution Grounding & 172 & 1 & 27 & 86\% & 86,5\% & 100\% & 72\% & 83,72\% \\
C -- HoarePrompt & 117 & 81 & 2 & 58,5\% & 99\% & 55,56\% & 85\% & 67,19\% \\
D -- LLM-as-a-Judge & 135 & 50 & 15 & 67,5\% & 92,5\% & 66,06\% & 72\% & 68,9\% \\
E -- Neuro-Symbolic & 16 & 0 & 184 & 8\% & 8\% & 100\% & 16\% & 27,59\% \\
F -- Agentic Multi-Step & 159 & 1 & 40 & 79,5\% & 80\% & 98,78\% & 81\% & 89,01\% \\
G -- Self-Consistency & 101 & 48 & 51 & 50,5\% & 74,5\% & 59,82\% & 67\% & 63,21\% \\
H -- Self-Refine & 179 & 21 & 0 & 89,5\% & 100\% & 86,24\% & 94\% & 89,95\% \\
I -- LLM-Guided Property-Based Testing & 100 & 0 & 100 & 50\% & 50\% & -- & 0\% & 0\% \\
J -- Deterministic Fuzzing & 100 & 0 & 100 & 50\% & 50\% & -- & 0\% & 0\% \\
K -- Mutation Testing & 91 & 60 & 49 & 45,5\% & 75,5\% & 34,48\% & 10\% & 15,5\% \\
L -- Symbolic Execution & 181 & 5 & 14 & 90,5\% & 93\% & 100\% & 81\% & 89,5\% \\
M -- Hoare Triple bằng ngôn ngữ tự nhiên & 109 & 66 & 25 & 54,5\% & 87,5\% & 59,2\% & 74\% & 65,78\% \\
N -- Invariant Synthesis & 20 & 0 & 180 & 10\% & 10\% & 100\% & 20\% & 33,33\% \\
\bottomrule
\end{tabular}
\end{table}

Trong Bảng \ref{tab:benchmark-overall-expanded}, \emph{Sai} đếm các phán quyết xác định nhưng sai nhãn. \emph{Chưa kết luận} đếm các mẫu chưa có phán quyết xác định. Precision, Recall và F1 lấy \texttt{VERIFIED} làm lớp dương. Precision là tỷ lệ đúng trong số phán quyết \texttt{VERIFIED}, còn Recall tính trên toàn bộ 100 mẫu có nhãn \texttt{VERIFIED}. Do đó, mẫu chưa kết luận có nhãn \texttt{VERIFIED} làm giảm Recall. Dấu gạch ngang ở cột Precision cho biết phương pháp chưa đưa ra phán quyết \texttt{VERIFIED} nào.

Hai chế độ Proposed đều đạt 200/200.

\section{Phân tích theo các câu hỏi nghiên cứu}
\label{sec:research_question_results}

\subsection{Khẳng định do LLM đề xuất và tính khả đạt (RQ1--RQ3)}

Proposed với MiniCPM ghi nhận 774 khẳng định bổ sung. Trong đó có 377 phán quyết \texttt{VERIFIED}, 206 \texttt{COUNTEREXAMPLE}, 144 \texttt{UNKNOWN\_TIMEOUT}, 28 \texttt{UNSUPPORTED} và 19 \texttt{TRANSLATION\_ERROR}. Các số đếm mô tả đầu ra của quy trình kiểm tra. Độ chính xác của các khẳng định bổ sung và tỷ lệ phát hiện lỗi của cổng neo cần một tập nhãn riêng ở cấp khẳng định. Bộ 200 mẫu chỉ có hai lớp nhãn chuẩn, nên chưa đo được độ nhạy của cơ chế phát hiện nhánh chết ở RQ3.

\subsection{Tái hiện phản ví dụ và phản hồi hồi cứu (RQ4)}
\label{sec:crvr_deep_dive}

Trên các khẳng định gốc, mỗi chế độ Proposed báo đúng 100 mẫu \texttt{COUNTEREXAMPLE}; cả 100 phản ví dụ trong mỗi chế độ đều được CPython xác nhận. CRVR trên nhóm khẳng định gốc này là 100\%. Phán quyết \texttt{COUNTEREXAMPLE} của C và H là đầu ra của mô hình cho bài toán phân loại. Hai lượt này không có kết quả tái hiện phản ví dụ. Thực nghiệm chưa định lượng tác động riêng của vòng phản hồi hồi cứu.

\subsection{Độ chính xác và tài nguyên vận hành (RQ5)}
\label{sec:cost_latency_tradeoffs}

Proposed Offline và Proposed với MiniCPM có cùng độ chính xác trên các khẳng định gốc. Ở phép loại bỏ cơ chế bất biến trong chế độ offline, hệ thống tìm được 100 phản ví dụ và để lại 100 mẫu chưa kết luận. Kết quả này cho thấy đóng góp của cơ chế bất biến vào độ phủ trên tập vòng lặp hiện tại.

\begin{table}[htbp]
\centering
\small
\caption{Tài nguyên ghi nhận trong các lượt đại diện trên 200 mẫu}
\label{tab:benchmark-resources}
\begin{tabular}{@{}lrrr@{}}
\toprule
\textbf{Phương pháp} & \textbf{Token vào} & \textbf{Token ra} & \textbf{Tổng giây} \\
\midrule
Proposed Offline & 0 & 0 & 97,89 \\
Proposed + MiniCPM & 144.127 & 48.536 & 2.659,37 \\
L -- Symbolic Execution & 0 & 0 & 375,60 \\
C -- HoarePrompt & 4.231.463 & 1.742.954 & 50.683,18 \\
H -- Self-Refine & 860.530 & 273.857 & 45.134,40 \\
\bottomrule
\end{tabular}
\end{table}

Cột thời gian là tổng thời lượng đã ghi từ toàn bộ quá trình tạo kết quả cuối cùng. Thời lượng có thể bao gồm thời gian chờ API và chờ theo giới hạn tốc độ gọi. Các phương pháp có tốc độ gọi và ngân sách thực thi khác nhau.

\section{Phạm vi kết luận và khả năng tái lập}
\label{sec:threats_to_validity}

\begin{itemize}
    \item \textbf{Phạm vi dữ liệu.} Tất cả mẫu thuộc họ Code2Inv và dùng ngữ nghĩa số nguyên của tập con Python. Một số mẫu cùng họ có cấu trúc tương tự. Khả năng tổng quát hóa cần được đo trên bộ chương trình độc lập, đa dạng hơn.
    \item \textbf{Phạm vi phán quyết.} Replay xác nhận từng đầu vào phản ví dụ cụ thể. Các phép đo trên khẳng định gốc giới hạn ở tính đúng từng phần tại vị trí câu lệnh.
\end{itemize}

Bảng được tạo từ \path{results/benchmark_200.csv}. Danh sách định danh mẫu nằm trong \path{results/checkpoints/results_minicpm_original_asserts_200/program_ids.json}. Tệp CSV chỉ đến các bản ghi nguồn của kết quả cuối cùng cho mỗi phương pháp.
